\documentclass[11pt]{amsart}
\usepackage[margin=1.1in]{geometry}
\usepackage{amsmath,amssymb,amsthm}
\usepackage{booktabs}
\usepackage[colorlinks=true,linkcolor=blue,citecolor=blue,urlcolor=blue]{hyperref}

\newtheorem{theorem}{Theorem}[section]
\newtheorem{proposition}[theorem]{Proposition}
\newtheorem{lemma}[theorem]{Lemma}
\newtheorem{corollary}[theorem]{Corollary}
\theoremstyle{definition}
\newtheorem{definition}[theorem]{Definition}
\newtheorem{hypothesis}[theorem]{Hypothesis}
\theoremstyle{remark}
\newtheorem{remark}[theorem]{Remark}

\newcommand{\lbl}[1]{\textup{\textsc{[#1]}}}

\title[The parity barrier and twin primes]{The parity barrier and twin primes: what is known and what is missing}
\author{Jerry Kim}
\thanks{Prepared with the assistance of Claude (Anthropic). Scripts and notes are in the repository \url{https://github.com/jerryckim16/twin-primes}; the companion papers are listed at \url{https://jerryckim16.github.io/twin-primes/}.}
\date{Version 0.1 (draft), 15 November 2026}

\begin{document}

\begin{abstract}
This is a survey-with-checks, not a research announcement: the twin prime
conjecture is open and nothing here proves it or comes close. We state, with
exact quotations checked in the sources, the parity obstruction that stops
sieve methods of the Maynard--Tao type at $H_1\le 6$ even under the
generalised Elliott--Halberstam conjecture, and the Friedlander--Iwaniec
asymptotic sieve, the one sieve framework with a bilinear axiom that breaks
parity. For the twin-prime sequence $a_n=\Lambda(n+2)$ we check that every
hypothesis of the asymptotic sieve except the Type~I and bilinear ones is satisfied, so that the
Hardy--Littlewood asymptotic would follow from a Type~I bound at level
$x^{2/3+}$ and one bilinear bound; we show that the bilinear bound has no
cancellation at all for the parity-twisted sequence, and that Cauchy--Schwarz
only moves the parity problem into an off-diagonal correlation. We also note
that the input of the parity heuristic, smallness of Liouville-twisted
Type~I sums, is a theorem in the logarithmically averaged, bounded-level
regime (a direct consequence of Tao's two-point theorem), and we give a small
numerical illustration. Every statement carries a label saying how it was
checked.
\end{abstract}

\maketitle

\section{Introduction}

Let $H_1=\liminf_n (p_{n+1}-p_n)$. Zhang, Maynard and Tao, and Polymath8b
proved $H_1\le 246$ \cite{PM8b}; several 2026 preprints claim smaller bounds,
and the companion papers of this project re-check the numerical parts of those
claims. The twin prime conjecture is the statement $H_1=2$. It is open.

Sieve methods that control primes only through \emph{linear} information
(how primes and related functions are distributed in arithmetic progressions)
cannot reach $H_1=2$, and in the Maynard--Tao framework cannot go below
$H_1\le 6$, even if the strongest conjectured linear input is granted. This
is the parity barrier. The aim of this paper is narrow: to state precisely
what the barrier is, which known results get around it and why their
mechanism does not obviously transfer to twin primes, and exactly which
estimate is missing. It grows run by run (see the changelog), and new
material enters only after it has been read in its source or proved in full.

\subsection*{Labels}
\lbl{Literature}: read in the cited source (version and location given).
\lbl{Proved here}: full proof given here; ``possibly folklore'' unless a
literature search says otherwise. \lbl{Heuristic}: numerical or informal;
never used as an input. \lbl{Open}: not known.

\section{The parity obstruction for sieve methods}\label{sec:parity}

\subsection{Selberg's example}
\lbl{Literature} \cite[pp.~1044--1045]{FI98}. Let $\lambda(n)=(-1)^{\Omega(n)}$
be the Liouville function and $a_n=\tfrac12(1+\lambda(n))$, the indicator of
integers with an even number of prime factors; then $a_p=0$ for every prime
$p$. Friedlander and Iwaniec write that this sequence ``satisfies (R) with
$D=x^{1-\varepsilon}$ but fails to satisfy (B)'', where (R) is the Type~I
remainder hypothesis and (B) the bilinear hypothesis of
Section~\ref{sec:fi}. In the classical set-up, (R) alone ``cannot capture
the primes, even for $D(x)=x^{1-\varepsilon}$'' \cite[p.~1044]{FI98}: the
sequences $\tfrac12(1+\lambda(n))$ and $\tfrac12(1-\lambda(n))$ have the same
Type~I data up to negligible error, but only the second contains primes.

\subsection{The obstruction for $H_1$}
\lbl{Literature} \cite[\S8, pp.~68--71]{PM8b}. Polymath8b argue, ``somewhat
informal and heuristic in nature'', that sieve methods cannot prove anything
stronger than $H_1\le 6$, ``even on the assumption of strong distributional
conjectures such as the generalized Elliott-Halberstam conjecture'', and
``even if one uses sieves other than the Selberg sieve''. The structure of the
argument is the following.
\begin{enumerate}
\item A sieve proof of $\mathrm{DHL}$-type statements produces a lower bound
  $\sum_n \nu(n)1_A(n)>0$ for a non-negative weight $\nu$ supported on
  $[x,2x]$, derived only from upper bounds for the discrepancies
  \cite[(142)]{PM8b}
  \[
  \left|\sum_{\substack{x\le n\le 2x\\ n\equiv a\ (q)}} f(n+h)
   -\frac{1}{\varphi(q)}\sum_{\substack{x\le n\le 2x\\ (n+h,q)=1}} f(n+h)\right|
  \]
  for $h$ in the tuple, $q\le x^{1-\varepsilon}$ and $f\in\{1,\Lambda\}$ or
  $f=\alpha\star\beta$ as in GEH, together with the main terms
  \cite[(143)]{PM8b}.
\item The same argument, run with a further weight $\omega\ge0$ inserted in
  every sum over $n$, would give $\sum_n\nu(n)1_A(n)\omega(n)>0$, provided
  the weighted discrepancies and main terms \cite[(144)--(145)]{PM8b} are
  controlled in the same way.
\item For $A_1=\{n: n,n+2\text{ prime}\}\cup\{n:n+2,n+6\text{ prime}\}$ (the
  set needed for $H_1\le 4$) take
  $\omega(n)=(1-\lambda(n)\lambda(n+2))(1-\lambda(n+2)\lambda(n+6))$. Since
  $\lambda(p)\lambda(p')=1$ for primes $p,p'$, $\omega$ vanishes on $A_1$, so
  the weighted lower bound is false \cite[(148)]{PM8b}.
\item The ``M\"obius randomness law'' predicts that the extra terms created by
  $\omega$, sums such as $\sum_{n\equiv a\,(q)} f(n)\lambda(n+2)\lambda(n+6)$,
  are negligible; Polymath8b ``conclude (heuristically, at least) that all the
  bounds that are believed to hold for (142), (143) should also hold (up to
  minor changes in the implied constants) for (144), (145)''
  \cite[p.~70]{PM8b}. Note $(\alpha\star\beta)\lambda=(\alpha\lambda)\star
  (\beta\lambda)$, so the twist preserves the GEH class.
\end{enumerate}
For twin primes themselves the simpler weight $\omega(n)=1-\lambda(n)
\lambda(n+2)$ suffices, ``and the argument is essentially due to Selberg''
\cite[p.~70]{PM8b}.

\begin{remark}[What kind of statement this is]\label{rem:meta}
The obstruction is not a theorem about all proofs. It says: any argument that
uses only the linear inputs (142)--(143), in a way that is insensitive to
inserting a non-negative weight, would also prove a false statement, as long
as the twisted inputs (144)--(145) hold. Two things are heuristic: that the
twisted inputs hold at the level needed (this is a family of Chowla-type
conjectures in progressions), and that ``sieve-theoretic'' arguments are
weight-insensitive in this sense. Proposition~\ref{prop:tao-weight} below
shows that the first is a theorem in a restricted (logarithmically averaged,
bounded-level) regime.
\end{remark}

\subsection{The twisted Type~I input in the log-averaged regime}

\lbl{Literature} \cite[Theorem~1.2]{Tao16} (arXiv:1509.05422): let
$a_1,a_2\in\mathbb N$, $b_1,b_2\in\mathbb Z$ with $a_1b_2-a_2b_1\ne0$, and
$1\le\omega(x)\le x$ with $\omega(x)\to\infty$; then
$\sum_{x/\omega(x)<n\le x}\lambda(a_1n+b_1)\lambda(a_2n+b_2)/n=o(\log\omega(x))$.
Tao notes that such estimates ``are well known to be subject to the parity
problem obstruction'' and that the proof gets around it with a new bilinear
estimate that uses the multiplicativity $\lambda(pn)=-\lambda(n)$ at small
primes; and that the method does ``not appear to have any bearing as yet on
twin prime-type sums'' \cite[pp.~3--5]{Tao16}.

\begin{proposition}[\lbl{Proved here}; possibly folklore]\label{prop:tao-weight}
Let $w:\mathbb N\to\mathbb R$ be periodic with period $q$. Then
\[
\sum_{n\le x}\frac{w(n)\lambda(n)\lambda(n+2)}{n}=o(\log x)\qquad(x\to\infty).
\]
In particular, if $R\ge1$ is fixed, $(\lambda_d)_{d\le R}$ are real with
$\lambda_1=1$, and $w(n)=\left(\sum_{d\mid n(n+2),\,d\le R}\lambda_d\right)^2$,
then with $c_w:=\lim_{x\to\infty}(\log x)^{-1}\sum_{n\le x}w(n)/n>0$,
\[
\sum_{\substack{n\le x\\ \lambda(n)\lambda(n+2)=\pm1}}\frac{w(n)}{n}
=\left(\tfrac12+o(1)\right)c_w\log x
\]
for both signs. The same holds with $R=R(x)\to\infty$ sufficiently slowly.
\end{proposition}

\begin{proof}
Write $w$ as a linear combination of the indicators of the residue classes
$c\bmod q$, $1\le c\le q$. For one class, $n=qm+c$ with $m\ge0$; the term
$m=0$ contributes $O(1)$, and for $m\ge1$, $1/n=1/(qm)+O(1/m^2)$. Hence
\[
\sum_{\substack{n\le x\\ n\equiv c\ (q)}}\frac{\lambda(n)\lambda(n+2)}{n}
=\frac1q\sum_{1\le m\le (x-c)/q}\frac{\lambda(qm+c)\lambda(qm+c+2)}{m}+O(1).
\]
Apply \cite[Theorem~1.2]{Tao16} with $a_1=a_2=q$, $b_1=c$, $b_2=c+2$
(so $a_1b_2-a_2b_1=2q\ne0$) and $\omega(y)=y$ at $y=(x-c)/q$: the sum over
$m$ is $o(\log x)$. Summing over the $q$ classes gives the first claim.

For the second, $w$ is periodic modulo $q=\prod_{d\le R}d$ (indeed modulo the
least common multiple of the $d$), because $1_{d\mid n(n+2)}$ depends only on
$n\bmod d$. A periodic $w\ge0$ has $\sum_{n\le x}w(n)/n=c_w\log x+O_w(1)$ with
$c_w$ the mean of $w$ over a period. If $n$ is odd and $n\not\equiv0,-2
\pmod p$ for every odd prime $p\le R$ then the only $d\le R$ dividing $n(n+2)$
is $d=1$, so $w(n)=\lambda_1^2=1$; such $n$ form a non-empty union of classes
modulo $\prod_{p\le R}p$ (choose $n\equiv1\pmod 2$, $n\equiv 2\pmod 3$ and,
for $p\ge5$, any $n\not\equiv0,-2$), hence $c_w>0$. Since
$1_{\lambda(n)\lambda(n+2)=\pm1}=\tfrac12(1\pm\lambda(n)\lambda(n+2))$ for
$n\ge1$, the first claim gives the second. Finally, for each fixed $R$ the
$o(1)$ above tends to $0$; a standard diagonal choice of $R(x)\to\infty$
(take $R(x)=R_j$ on $[x_j,x_{j+1})$ with $x_j$ so large that the error for
all $R\le R_j$ is below $1/j$, and with the weights chosen from a fixed
sequence) gives the last statement.
\end{proof}

\begin{remark}[What is missing]\label{rem:tao-missing}
Sieve weights for twin primes need level $R=x^{\theta}$ with $\theta>0$
fixed. Proposition~\ref{prop:tao-weight} needs $R$ fixed or growing as slowly
as the (ineffective, see \cite[Remark~1.4 and footnote~2]{Tao16}) dependence
of Tao's theorem on $a_1,b_1,b_2$ allows. Uniformity of
\cite[Theorem~1.2]{Tao16} over moduli $q\le x^{\theta}$ is \lbl{Open} as far
as we know. Even with it, Proposition~\ref{prop:tao-weight} would only
confirm the input of the obstruction: it says that the sieve weight cannot
tell $\lambda(n)\lambda(n+2)=+1$ from $-1$, which is exactly why it cannot
see twin primes. It does not give a lower bound for twin primes.
\end{remark}

\section{The asymptotic sieve and the twin-prime sequence}\label{sec:fi}

\subsection{The theorem}
\lbl{Literature} \cite{FI98} (arXiv:math/9811186v1, journal pagination). Let
$a_n\ge0$, $A(x)=\sum_{n\le x}a_n$, $A_d(x)=\sum_{n\le x,\,d\mid n}a_n=g(d)A(x)
+r_d(x)$ with $g$ multiplicative. The hypotheses are: (1.4)
$A(x)\gg A(\sqrt x)(\log x)^2$; (1.6) $A_d(x)\ll d^{-1}\tau(d)^8A(x)$ for
$d\le x^{1/3}$; (1.8) $0\le g(p)<1$, $g(p)\ll p^{-1}$; (1.9)
$\sum_{p\le y}g(p)=\log\log y+c+O((\log y)^{-10})$; the Type~I hypothesis
\[
\text{(R)}\qquad \sum_{d\le D}\mu^2(d)|r_d(t)|\le A(x)(\log x)^{-2^{22}}
\quad(t\le x),\qquad \text{(R1)}\quad x^{2/3}<D(x)<x;
\]
and the bilinear hypothesis
\[
\text{(B)}\qquad \sum_m\left|\sum_{\substack{N<n\le2N\\ mn\le x}}
\gamma(n,C)\mu(mn)a_{mn}\right|\le A(x)(\log x)^{-2^{22}},
\qquad \gamma(n,C)=\sum_{d\mid n,\,d\le C}\mu(d),
\]
for every $N$ with $\Delta^{-1}\sqrt D<N<\delta^{-1}\sqrt x$ and every
$1\le C\le xD^{-1}$. Theorem~1 \cite[p.~1044]{FI98} (for $a_n$ supported on
squarefree $n$) gives
$\sum_{p\le x}a_p\log p=HA(x)\{1+O(\log\delta/\log\Delta)\}$ with
$H=\prod_p(1-g(p))(1-1/p)^{-1}$. Theorem~2 \cite[\S9, p.~1059]{FI98} drops
the squarefree condition, adding (9.1) $0\le g(p^2)\le g(p)$, $g(p^2)\ll
p^{-2}$, (9.2) $\sum_{n\le x}a_n^2\le x^{-2/3}A(x)^2$, and replacing (R) by
the cubefree version (R$_3$) $\sum^{3}_{d<DL^2}|r_d(t)|\le A(x)L^{-2}$,
$L=(\log x)^{2^{24}}$. Friedlander and Iwaniec remark that the lower limit
$N>\Delta^{-1}\sqrt D$ in (B) ``is essential; indeed by narrowing this
slightly to $N>\sqrt D$ we would not be able to break the parity problem''
\cite[p.~1045]{FI98}, and, on twin primes \cite[p.~1046]{FI98}: ``The
stumbling block is that in this case we have no idea how to prove that the
relevant sequence satifies the condition (B). Even proving (R) to such a high
level is currently beyond reach.''

\subsection{Specialisation to $a_n=\Lambda(n+2)$}

\begin{proposition}[\lbl{Proved here}, modulo the Brun--Titchmarsh inequality]
\label{prop:fi-twin}
For $a_n=\Lambda(n+2)$ the hypotheses (1.4), (1.6), (1.8), (1.9), (9.1),
(9.2) of \cite[Theorem~2]{FI98} hold with $g(2^j)=0$ and $g(p^j)=1/\varphi(p^j)$
for odd $p$, and $H=2C_2$ with $C_2=\prod_{p\ge3}p(p-2)/(p-1)^2$. Hence, if
\textup{(R$_3$)} holds with some $D>x^{2/3}$ and \textup{(B)} holds for this
sequence in the stated ranges, then
$\sum_{p\le x}\Lambda(p+2)\log p=2C_2x\{1+O(\log\delta/\log\Delta)\}$,
which is the Hardy--Littlewood twin-prime asymptotic.
\end{proposition}

\begin{proof}
For odd $d$, $A_d(t)=\sum_{m\le t+2,\ m\equiv2\,(d)}\Lambda(m)+O(\log t)
=\psi(t+2;d,2)+O(\log t)$, and $(2,d)=1$, so the natural density is
$1/\varphi(d)$. For even $d$, $d\mid n$ and $\Lambda(n+2)\ne0$ force
$n+2$ to be a power of $2$, so $A_d(t)=O((\log t)^2)$, and we take $g(d)=0$.
So $g(p^j)=1/\varphi(p^j)$ for odd $p$ and $g(2^j)=0$. Then $0\le g(p)\le
\tfrac12<1$ and $g(p)\ll1/p$, which is (1.8); (1.9) follows from Mertens'
theorem since $\sum_{p\le y}1/(p-1)=\log\log y+c'+O(1/\log y)$ with the prime
number theorem error (much better than $(\log y)^{-10}$); (9.1) holds since
$g(p^2)=1/(p(p-1))\le 1/(p-1)$ and $\ll p^{-2}$. By the prime number theorem
$A(x)\sim x$, giving (1.4), and $\sum_{n\le x}a_n^2\le x\log^2(x+2)\ll
x^{4/3}\ll x^{-2/3}A(x)^2$, giving (9.2). For (1.6), Brun--Titchmarsh gives
$A_d(x)\ll x/(\varphi(d)\log(x/d))\ll x/\varphi(d)$ for odd $d\le x^{1/3}$, and
$d/\varphi(d)\le 2^{\omega(d)}\le\tau(d)$. Finally
$H=\prod_p(1-g(p))(1-1/p)^{-1}=2\prod_{p\ge3}\frac{p-2}{p-1}\cdot\frac{p}{p-1}
=2C_2$.
\end{proof}

So two inputs are missing. (R$_3$) at a level $D>x^{2/3}$ is a fixed-residue
Elliott--Halberstam bound with absolute values; the best known level of this
form is $x^{1/2-\varepsilon}$ (Bombieri--Vinogradov, as stated in
\cite[Theorem~2.3]{PM8b}). (B) for $\Lambda(n+2)$ is the bilinear estimate
discussed next. Whether there is progress after 2014 on absolute-value,
fixed-residue level beyond $x^{1/2}$ has not yet been checked (Section~\ref{sec:open}).

\subsection{The bilinear hypothesis carries the parity information}

\begin{proposition}[\lbl{Proved here}]\label{prop:B-parity}
Let $a_n=\Lambda(n+2)$ and $a^{+}_n=\Lambda(n+2)\tfrac12(1+\lambda(n))$.
\begin{enumerate}
\item With $C=1$, the inner sum of \textup{(B)} for $a$ is
  $\sum_n\mu^2(mn)\lambda(mn)\Lambda(mn+2)$: \textup{(B)} at $C=1$ is a bilinear
  estimate for $k\mapsto\lambda(k)\Lambda(k+2)$.
\item With $C=1$, every term $\mu(mn)a^+_{mn}$ is $\ge0$, so the left side of
  \textup{(B)} for $a^+$ equals its trivial bound
  $\sum_{m}\sum_{N<n\le2N,\,mn\le x}\mu^2(mn)a^+_{mn}$.
\item Bound the left side of \textup{(B)} for $a$ at $C=1$ by Cauchy--Schwarz in
  $m$. The diagonal $n_1=n_2$ is acceptable, and the following off-diagonal bound
  suffices for \textup{(B)} at $C=1$:
  \[
  \text{(OD)}\qquad\left|\sum_{\substack{n_1\ne n_2\\ N<n_i\le 2N}}\mu(n_1)\mu(n_2)
  \sum_{\substack{m\le x/\max(n_1,n_2)\\(m,n_1n_2)=1}}\mu^2(m)
  \Lambda(mn_1+2)\Lambda(mn_2+2)\right|\le \tfrac14 Nx(\log x)^{-2^{23}}.
  \]
  For $a^+$ in place of $a$ every term of the corresponding sum is $\ge0$.
\end{enumerate}
\end{proposition}

\begin{proof}
(1) $\gamma(n,1)=1$, and $\mu(k)\ne0$ implies $\mu(k)=\lambda(k)$.
(2) If $\mu(k)\ne0$ then $\mu(k)\cdot\tfrac12(1+\lambda(k))=\tfrac12(\mu(k)+1)
\in\{0,1\}$, and $\Lambda\ge0$.
(3) Put $S_m=\sum_{N<n\le2N,\,mn\le x}\mu(mn)a_{mn}$; only $m\le x/N$ occur.
By Cauchy--Schwarz, $(\sum_m|S_m|)^2\le (x/N)\sum_m|S_m|^2$, and for
$(m,n)=1$, $\mu(mn)=\mu(m)\mu(n)$, so
$\sum_m|S_m|^2=\sum_{n_1,n_2}\mu(n_1)\mu(n_2)\sum_{m}\mu^2(m)
a_{mn_1}a_{mn_2}$ with $m\le x/\max(n_1,n_2)$, $(m,n_1n_2)=1$ and $n_i$
squarefree (terms with $\mu(mn_i)=0$ vanish). The diagonal is
$\sum_n\sum_{m\le x/n}\Lambda(mn+2)^2\ll\sum_{n\le2N}(x/\varphi(n))\log x\ll
x(\log x)^2$ by Brun--Titchmarsh (for odd $n$; even $n$ contribute only
prime powers of $2$), so its share of $(\sum_m|S_m|)^2$ is $\ll
(x^2/N)(\log x)^2$, which is $\le\tfrac14 x^2(\log x)^{-2^{23}}$ as soon as
$N\ge(\log x)^{2^{23}+3}$ and $x$ is large, true in the range of (B). With (OD)
the off-diagonal share is at most $(x/N)\cdot\tfrac14Nx(\log x)^{-2^{23}}$, so
$(\sum_m|S_m|)^2\le\tfrac12x^2(\log x)^{-2^{23}}\le A(x)^2(\log x)^{-2^{23}}$
for large $x$, since $A(x)\sim x$. For $a^+$: if $a^+_{mn_i}\ne0$ with $mn_i$
squarefree then $\lambda(mn_i)=\mu(m)\mu(n_i)=1$, so $\mu(n_i)=\mu(m)$ and
$\mu(n_1)\mu(n_2)=1$; all terms are $\ge0$.
\end{proof}

\begin{remark}
Proposition~\ref{prop:B-parity}(2) says (B) separates $a$ from $a^+$; by the
heuristic of Section~\ref{sec:parity}, Type~I data do not
\cite[p.~70]{PM8b}. That (B) actually fails for $a^+$ needs the trivial bound to
be $\gg A(x)(\log x)^{-2^{22}}$, i.e.\ enough primes $q$ with
$\lambda(q-2)=1$ and $q-2$ having a squarefree divisor in $(N,2N]$; this is
expected but we have not proved it. Polymath8b say of such bilinear sums for
$\Lambda(dm+2)$: ``we do not know of any plausible way to control such
bilinear expressions'' \cite[p.~71]{PM8b}.
\end{remark}

\section{A numerical illustration}\label{sec:num}

\lbl{Heuristic.} For odd $n\le X$ let
\[
w(n)=\left(\sum_{\substack{d\mid n(n+2),\ d\le R\\ d\text{ odd squarefree}}}\mu(d)
\left(\frac{\log(R/d)}{\log R}\right)^2\right)^2,\qquad R=X^{\theta}
\]
(a GPY-type weight for the pair $\{0,2\}$). The script (file
\texttt{parity\_weighted\_liouville.py} in \texttt{code/explore/})
computes $B/A=\sum w(n)\lambda(n)\lambda(n+2)/\sum w(n)$ and the share
$T/A$ of the weight on twin-prime pairs:

\begin{center}
\begin{tabular}{rrrrr}
\toprule
$X$ & $\theta$ & $R$ & $B/A$ & $T/A$\\
\midrule
$10^7$ & $0.20$ & $25.1$ & $3.6\cdot10^{-4}$ & $0.0325$\\
$10^7$ & $0.25$ & $56.2$ & $6.0\cdot10^{-4}$ & $0.0490$\\
$10^8$ & $0.20$ & $39.8$ & $4.2\cdot10^{-5}$ & $0.0310$\\
$10^8$ & $0.25$ & $100$ & $4.3\cdot10^{-5}$ & $0.0466$\\
\bottomrule
\end{tabular}
\end{center}

So the two parity-twisted measures $w(1\pm\lambda(n)\lambda(n+2))/2$ carry
essentially equal mass, and (at $X=10^6$, printed by the script) the same
residue-class distributions modulo $3,5,7$ to about three decimals, while all
the twin-prime pairs, about $3$--$5\%$ of the weight, lie in the ``$+$''
measure. This is Selberg's phenomenon made visible; it proves nothing.

\section{Open questions}\label{sec:open}

\begin{enumerate}
\item (\lbl{Open}) Is \cite[Theorem~1.2]{Tao16} uniform in the modulus: is
  $\sum_{n\le x,\,n\equiv c\,(q)}\lambda(n)\lambda(n+2)/n=o(\log x/q)$
  uniformly for $q\le x^{\theta}$, on average over $q$? This is the twisted
  Type~I input of Section~\ref{sec:parity} in log-averaged form; even a positive
  answer would confirm the obstruction rather than break it.
\item Is there post-2014 progress on absolute-value, fixed-residue
  equidistribution of primes beyond level $x^{1/2}$, so that (R$_3$) at
  $x^{2/3+}$ is not a second, independent barrier? (Not yet checked.)
\item Is there a sequence closely tied to twin primes (for instance with a
  norm-form structure, as in Friedlander--Iwaniec's $a^2+b^4$) for which a
  (B)-type estimate is provable? Tao's blog post on the parity problem
  (2007) observes that the known parity-breaking methods ``all seem to
  require that the form be representable as a norm over some number field''
  (to be re-read before it is used).
\end{enumerate}

\subsection*{Plan for later versions}
Each run adds one bounded step: Friedlander--Iwaniec $x^2+y^4$ and
Heath-Brown $x^3+2y^3$ (statements and the structure each uses);
Matom\"aki--Radziwi\l\l(--Tao) and Tao--Ter\"av\"ainen; Sawin--Shusterman in
$\mathbb F_q[T]$ and the step with no integer analogue; conditional
statements (which hypotheses imply twin primes, which are known not to
suffice). Nothing enters before it is read in the source.

\begin{thebibliography}{9}
\bibitem{FI98} J.~Friedlander and H.~Iwaniec, \emph{Asymptotic sieve for primes},
  Ann.\ of Math.\ \textbf{148} (1998), 1041--1065; arXiv:math/9811186v1.
\bibitem{PM8b} D.~H.~J.~Polymath, \emph{Variants of the Selberg sieve, and bounded
  intervals containing many primes}, arXiv:1407.4897v4 (22 Dec 2014) (cited by its page numbers, of 80).
\bibitem{Tao16} T.~Tao, \emph{The logarithmically averaged Chowla and Elliott
  conjectures for two-point correlations}, arXiv:1509.05422 (journal version not
  consulted).
\end{thebibliography}

\section*{Changelog}
\begin{itemize}
\item[0.1] (15 Nov 2026) First version: Sections 2--3 from the sources read
  on 14--15 Nov (SHA-256 of the files in \texttt{notes/parity-programme.md});
  Proposition~\ref{prop:fi-twin} and Proposition~\ref{prop:B-parity} from
  the 14 Nov route card (\texttt{notes/parity-attempts/}), re-derived;
  Proposition~\ref{prop:tao-weight} (parity attempt of 15 Nov); numerical
  illustration.
\end{itemize}

\end{document}
